\subsection{Moderated functions and measures}

PROPOSITION 5. --- Let A be a subset of T; the following properties are equivalent:

\begin{enumerate}
\def\labelenumi{\alph{enumi})}
\item
  The set A is contained in the union of a sequence of \(\mu\) -integrable open sets.
\item
  The set A is contained in the union of a sequence of \(\mu\) -integrable sets.
\item
  The set A is contained in the union of a sequence of compact sets and a \(\mu\) -negligible set.
\end{enumerate}

It is clear that each of the properties a) and c) implies b). Conversely, b) implies a) because every set of finite outer measure is contained in an integrable open set (Ch. IV, §1, No.~4, Prop. 19), and b) implies c) because every integrable set is the union of a sequence of compact sets and a negligible set (Ch. IV, §4, No.~6, Cor. 2 of Th. 4).

DEFINITION 2. --- A subset of T is said to be \(\mu\) -moderated if it satisfies the equivalent conditions of Proposition 5. A function defined on T, with values in a vector space or in \(\overline{R}\) , is said to be \(\mu\) -moderated if it is zero on the complement of a \(\mu\) -moderated subset of T. The measure \(\mu\) is said to be moderated if T is a \(\mu\) -moderated set.

If \(\mu\) is a moderated measure, then every function on T is \(\mu\) -moderated and every subset of T is \(\mu\) -moderated.

Remarks. --- 1) If \(\theta\) is a complex measure on T, one says that a function \(f\) is \(\theta\) -moderated (resp. that \(\theta\) is moderated) if \(f\) is \(|\theta|\) -moderated (resp. if \(|\theta|\) is moderated).

\begin{enumerate}
\def\labelenumi{\arabic{enumi})}
\setcounter{enumi}{1}
\item
  Every bounded measure is moderated; if \(\mathrm{T}\) is a countable union of compact sets, then every measure on \(\mathrm{T}\) is moderated.
\item
  Let \((f_n)\) be a sequence of \(\mu\) -moderated functions with values in \(\overline{\mathbf{R}}\) . For each \(n\) , let \(\mathrm{U}_n\) be an open set that is a countable union of open sets of finite outer measure, such that \(f_n\) is zero outside \(\mathrm{U}_n\) . The function \(s = \sum_{n \in \mathbf{N}} |f_n|\) is then zero outside \(\bigcup_{n\in \mathbf{N}}\mathrm{U}_n\) ; it is therefore \(\mu\) -moderated, and the same is true
\end{enumerate}

of all functions \(f\) such that \(|f| \leqslant s\) . This applies in particular to the functions \(\liminf_{n \to \infty} f_n\) , \(\limsup_{n \to \infty} f_n\) and \(\sum_{n \in \mathbf{N}} f_n\) (if the sum is defined).

\begin{enumerate}
\def\labelenumi{\arabic{enumi})}
\setcounter{enumi}{3}
\tightlist
\item
  A function equal almost everywhere to a moderated function is moderated.
\end{enumerate}

PROPOSITION 6. --- Let f be a positive numerical function defined on T that is \(\mu\) -measurable and \(\mu\) -moderated. Then there exists a sequence \((h_{n})_{n\in\mathbf{N}}\) of elements of \(\mathcal{F}_{+}(\mathrm{T})\) , with sum equal to f, having the following properties:

\begin{enumerate}
\def\labelenumi{\arabic{enumi})}
\item
  The function \(h_0\) is \(\mu\) -negligible.
\item
  For every \(n \geqslant 1\) , there exists a compact set \(K_{n}\) such that \(h_{n}\) is zero outside \(K_{n}\) , and such that the restriction of \(h_{n}\) to \(K_{n}\) is finite and continuous.
\end{enumerate}

Suppose that \(f\) is the sum of a sequence \((f_n)\) of positive measurable functions, each of which has the property in the statement; it is clear that \(f\) also has it. Set

\[
f _ {n} = \inf (f, n + 1) - \inf (f, n)
\]

for every \(n \in \mathbf{N}\) ; since \(f\) is equal to the sum of the sequence \((f_n)\) , it will thus suffice to establish the proposition assuming \(f\) to be moderated and bounded. Denote then by \(A\) the set of \(t \in T\) such that \(f(t) > 0\) ; \(A\) is measurable and moderated, therefore there exists a sequence \((A_n)\) of pairwise disjoint integrable sets such that \(A = \bigcup_{n} A_n\) . We are reduced

to proving the statement for the functions \(f\varphi_{\mathrm{A}_n}\) ; in other words, we may suppose \(f\) to be bounded and to be zero outside an integrable set I. But I is the union of a negligible set N and a sequence \((\mathbf{L}_n)\) of pairwise disjoint compact sets (Ch. IV, §4, No.~6, Cor. 2 of Th. 4). We are thus reduced to treating the case that \(f\) is bounded and is zero outside a compact set L.

Let \(\mathfrak{K}\) be the set of compact subsets \(K\) of \(T\) such that \(f|K\) is continuous; since \(\mathfrak{K}\) is \(\mu\) -dense (Ch. IV, §5, No.~10, Prop. 15), \(L\) is the union of a negligible set \(N\) and a sequence \((K_n)_{n \geqslant 1}\) of pairwise disjoint elements of \(\mathfrak{K}\) (Ch. IV, §5, No.~8, Def. 6). The functions \(h_0 = f\varphi_N\) , \(h_n = f\varphi_{K_n}\) for \(n \geqslant 1\) then satisfy the conditions of the statement.

The following proposition makes it possible to reduce the study of the upper integral to that of the essential upper integral.

PROPOSITION 7. --- Let \(f\) be an element of \(\mathcal{F}_{+}(\mathrm{T})\) .

\begin{enumerate}
\def\labelenumi{\arabic{enumi})}
\item
  If the function \(f\) is not \(\mu\) -moderated, then \(\mu^{*}(f) = +\infty\) .
\item
  If the function \(f\) is \(\mu\) -moderated, then \(\mu^{*}(f) = \mu^{\bullet}(f)\) .
\item
  If \(\mu^{\bullet}(f) < +\infty\) then there exists a \(\mu\) -moderated subset A, the union of a sequence of compact subsets of T, such that \(f = f\varphi_{\mathrm{A}}\) locally almost everywhere.
\end{enumerate}

The first assertion follows immediately from Lemma 1 of Ch. IV, §5, No.~6. To establish the second, denote by A a moderated subset such that f is zero outside A; A is the union of a negligible set \(A_{0}\) and a sequence \((\mathrm{A}_{n})_{n\geqslant1}\) of compact sets, which we may suppose to be increasing. The function f is then almost everywhere equal to the upper envelope of the functions \(f\varphi_{\mathrm{A}_{n}}\) ( \(n\geqslant1\) ), therefore (Ch. IV, §1, No.~3, Th. 3 and §2, No.~3, Prop. 6)

\[
\mu^ {*} (f) = \lim _ {n \rightarrow \infty} \mu^ {*} (f \varphi_ {\mathrm{A} _ {n}}) \leqslant \mu^ {\bullet} (f),
\]

whence the equality \(\mu^{*}(f)=\mu^{\bullet}(f)\) by virtue of the formula (1). Finally, suppose that \(\mu^{\bullet}(f)<+\infty\) ; there exists an increasing sequence \((\mathrm{A}_{n})\) of compact sets such that

\[
\mu^ {\bullet} (f) = \sup _ {n} \mu^ {*} (f \varphi_ {\mathrm{A} _ {n}}).
\]

Set \(A = \bigcup_{n} A_{n}\) ; the second member is equal to \(\mu^{*}(f\varphi_{\mathrm{A}})\) (Ch. IV, §1, No.~3, Th. 3), that is, to \(\mu^{\bullet}(f\varphi_{\mathrm{A}})\) (by Prop. 1, or by 2) above). Since \(\mu^{\bullet}(f) = \mu^{\bullet}(f\varphi_{\mathrm{A}}) + \mu^{\bullet}(f\varphi_{\mathbb{C}_{\mathrm{A}}})\) (Prop. 2), we have \(\mu^{\bullet}(f\varphi_{\mathbb{C}_{\mathrm{A}}}) = 0\) , from which 3) follows.

COROLLARY 1. --- For f to be negligible, it is necessary and sufficient that it be locally negligible and moderated.

COROLLARY 2. --- If \(\mu\) is a moderated measure (in particular if \(\mu\) is bounded, or if T is countable at infinity), then \(\mu^{*} = \mu^{\bullet}\) .

PROPOSITION 8. --- a) Let H be a set of functions ≥ 0, lower semicontinuous, directed for the relation ≤; then

\[
\mu^ {\bullet} \left(\sup _ {h \in \mathrm{H}} h\right) = \sup _ {h \in \mathrm{H}} \mu^ {\bullet} (h).
\]

\begin{enumerate}
\def\labelenumi{\alph{enumi})}
\setcounter{enumi}{1}
\tightlist
\item
  Let H be a set of functions \(\geqslant0\) , upper semi-continuous, directed for the relation \(\geqslant\) ; if there exists in H a function \(h_{0}\) such that \(\mu^{\bullet}(h_{0})<+\infty\) , then
\end{enumerate}

\[
\mu^ {\bullet} \Bigl (\inf _ {h \in \mathrm{H}} h \Bigr) = \inf _ {h \in \mathrm{H}} \mu^ {\bullet} (h).
\]

The assertion a) is, in view of Prop. 4, a repetition of Theorem 1 of Ch. IV, §1, No.~1. To establish b), set \(\eta = \inf_{h\in \mathrm{H}}h\) , and let \(a\) be a number \(>0\) . There exists a compact set K such that (Ch. IV, §4, No.~4, Cor. 1 of Prop. 5):

\[
\mu^ {\bullet} (h _ {0}) - a \leqslant \mu^ {*} (h _ {0} \varphi_ {\mathrm{K}}) = \mu (h _ {0} \varphi_ {\mathrm{K}}) \leqslant \mu^ {\bullet} (h _ {0}).
\]

The functions \(h\varphi_{K}\) , where h runs over H, form a set of upper semicontinuous functions, directed for the relation \(\geqslant\) , which contains an integrable function. Therefore (Ch. IV, §4, No.~4, Cor. 2 of Prop. 5):

\[
\mu^ {*} (\eta \varphi_ {\mathrm{K}}) = \inf _ {h \in \mathrm{H}} \mu^ {*} (h \varphi_ {\mathrm{K}}).
\]

But (Ch. IV, §4, No.~4, Cor. 1 of Prop. 5) \(\mu^{\bullet}(h_{0}\varphi_{\mathbb{C}_{\mathbf{K}}}) \leqslant a\) , whence \(\mu^{\bullet}(h\varphi_{\mathbb{C}_{\mathbf{K}}}) \leqslant a\) for every function \(h \in \mathrm{H}\) such that \(h \leqslant h_{0}\) . Therefore, finally:

\[
\mu^ {\bullet} (\eta) \geqslant \mu^ {*} (\eta \varphi_ {\mathrm{K}}) = \inf _ {h \in \mathrm{H}} \mu^ {*} (h \varphi_ {\mathrm{K}}) \geqslant \inf _ {h \in \mathrm{H}} \mu^ {\bullet} (h) - a.
\]

The inequality \(\mu^{\bullet}(\eta) \leqslant \inf_{h \in \mathrm{H}} \mu^{\bullet}(h)\) being obvious, and a being arbitrary, the proposition is established.

